% Template modulaire FR pour CV — personnaliser via commandes en tête
\documentclass[10pt, letterpaper]{article}

\usepackage[
    top=2cm,
    bottom=2cm,
    left=2cm,
    right=2cm
]{geometry}
\usepackage{titlesec}
\usepackage{tabularx}
\usepackage{enumitem}
\usepackage[dvipsnames]{xcolor}
\definecolor{primaryColor}{RGB}{0, 79, 144}
\usepackage{fontawesome5}
\usepackage{hyperref}
\hypersetup{colorlinks=true, urlcolor=primaryColor}
\usepackage[french]{babel}
\usepackage{datetime2}
\usepackage{lastpage}
\usepackage{fancyhdr}
\usepackage{lmodern}

% --- Personnalisation (modifiez ces commandes avant de compiler) ---
\newcommand{\cvname}{Charly-Romy TANGA}
\newcommand{\cvaddress}{Paris, France}
\newcommand{\cvmail}{charlyromytanga.cytech@gmail.com}
\newcommand{\cvphone}{+33 6 15 42 25 74}
\newcommand{\cvlinkedin}{charly-romy.tanga}
\newcommand{\cvgithub}{charlyromytanga}
\newcommand{\jobtype}{Type de recherche : Alternance / CDI / Stage}
\newcommand{\presentation}{Remplacez ce texte par votre présentation synthétique. 1–2 lignes maximum.}

% ----- PIED DE PAGE -----
\pagestyle{fancy}
\fancyhf{}
\fancyfoot[R]{\small\textit{\color{gray}Dernière mise à jour le \DTMtoday}}
\fancyfoot[L]{\small\textit{\color{gray}\cvname{} – Page \thepage{} sur \pageref*{LastPage}}}
\renewcommand{\headrulewidth}{0pt}
\renewcommand{\footrulewidth}{0pt}

% ----- FORMAT DES SECTIONS -----
\titleformat{\section}{\large\bfseries\color{primaryColor}}{}{0pt}{}
\titlespacing*{\section}{0pt}{0.2cm}{0.15cm}
\setlist[itemize]{leftmargin=0.5cm, itemsep=3pt}

\begin{document}

% ===== EN-TÊTE =====
% Header: utilise les commandes definies dans main.tex
\begin{center}
    {\Large \textbf{\cvname}}\\[3pt]
    {\small \cvaddress{} \quad$\cdot$\quad \href{mailto:\cvmail}{\cvmail} \quad$\cdot$\quad
    Tél : \cvphone}\\[3pt]
    \href{https://linkedin.com/in/\cvlinkedin}{\faLinkedinIn\ \cvlinkedin} \quad$\cdot$\quad
    \href{https://github.com/\cvgithub}{\faGithub\ \cvgithub}
\end{center}

\vspace{0.2cm}


% ===== TYPE DE RECHERCHE =====
\section*{Type de recherche}
% Modifiez la commande \jobtype dans main.tex ou remplacez le contenu ci-dessous.
\noindent\textbf{\jobtype}


% ===== PRÉSENTATION =====
\section*{Présentation}
% Remplacez la commande \presentation dans main.tex ou editez ci-dessous.
\noindent\presentation


% ===== EXPÉRIENCES =====
\section*{Expériences}
% Exemples d'entrées — remplacez par vos expériences réelles
\textbf{Alternant Data Analyst – NaTran France (Spécialiste national des données gazières)} \hfill \textit{Septembre 2023 – Présent}\\
\textit{Ingénierie de la données, Département Data}
\begin{itemize}
    \item Développement de tableaux de bord nationaux sous \textbf{Power BI}, intégrant des jeux de données hétérogènes et réduisant le temps de reporting de 40\%.
    \item Conception et optimisation de pipelines de traitement de données avec \textbf{Python (Pandas, PySpark)}, automatisant des tâches récurrentes.
\end{itemize}

% Ajoutez d'autres expériences au besoin


% ===== CERTIFICATIONS =====
\section*{Certifications professionnelles}
% Certifications — listez vos certifications ici
\begin{itemize}
    \item Microsoft Certified: Data Analyst Associate (en cours) – Power BI
    \item Bloomberg Market Concepts (BMC) (en cours)
    \item Certification AMF (en cours)
\end{itemize}


% ===== COMPÉTENCES =====
\section*{Compétences}
% Compétences et technologies
\textbf{Langages :} Python, Java, C, SQL\\
\textbf{Outils :} Power BI, Alteryx (préparation), NumPy, pandas, SciPy, scikit-learn, PySpark, Git, Docker\\
\textbf{Maths/Finance :} Probabilités, Statistiques, Processus stochastiques, Finance computationnelle


% ===== FORMATION =====
\section*{Formation}
% Formation — adaptez les éléments ci-dessous
\textbf{Master 2 en Mathématiques appliquées et Informatique – Spécialisation FinTech} \hfill \textit{CY Tech, Cergy – 2024–2026}\\
\textbf{Licence en Mathématiques et Informatique} \hfill \textit{CY Tech, Cergy – 2023–2024}


\end{document}
